\documentclass[11pt, reqno]{article}

\usepackage{amsmath, amsthm, amssymb}
\usepackage{enumitem}
\usepackage{tcolorbox}
\usepackage{hyperref}
\usepackage{tikz}
\usepackage{tikz-cd}
\usepackage{pgfplots}
\pgfplotsset{compat=1.18}
\usetikzlibrary{arrows.meta}
\usepackage{mathrsfs}
\usepackage{fancyhdr}
\usepackage[bottom=0.75in, top=1in, left=0.5in, right=0.5in]{geometry}
\usepackage{array}   % for \newcolumntype macro
\newcolumntype{L}{>{$}l<{$}}

\theoremstyle{plain}
\newtheorem*{theorem}{Theorem}
\newtheorem*{proposition}{Proposition}
\newtheorem{exercise}{Exercise}
\newtheorem*{lemma}{Lemma}
\newtheorem*{corollary}{Corollary}

\theoremstyle{definition}
\newtheorem*{definition}{Definition}
\newtheorem*{example}{Example}

\theoremstyle{remark}
\newtheorem*{remark}{Remark}

\renewcommand{\phi}{\varphi}
\renewcommand{\epsilon}{\varepsilon}
\renewcommand{\emptyset}{\varnothing}

\newcommand{\RR}{\mathbb{R}}
\newcommand{\ZZ}{\mathbb{Z}}
\newcommand{\NN}{\mathbb{N}}
\newcommand{\CC}{\mathbb{C}}
\newcommand{\QQ}{\mathbb{Q}}

\DeclareMathOperator{\ima}{\text{im}}

\begin{document}

\topmargin=-40pt
\rhead{Henry Woodburn}
\lhead{Math 656}
\renewcommand{\headrulewidth}{1pt}
\renewcommand{\headsep}{20pt}
\thispagestyle{fancy}

{\Large \bfseries \noindent Dirac Operators, Spin manifolds, and Positive Scalar Curvature}

\tableofcontents

\section{Introduction}

These are the notes for my final presentation in Functional Analysis II. The goal of this presentation will 
be to ``prove'' a theorem about certain types of manifolds which do not admit positive scalar curvature 
using the Atiyah-Singer index theorem. Along the way, we will define the Dirac Operator on a Riemannian Manifold,
investigate spin structures, and see how these two interact on a spin manifold. 

\section{Riemannian Geometry}

\begin{definition}
    Let $M$ be a smooth manifold, and $V$ a vector bundle on $M$. 
    A \textbf{Connection} on $V$ is an $\RR$-linear map $\nabla: C^\infty(V) \to \Omega(V) = C^\infty(T^*M \otimes V)$ which
    satisfies the Leibniz rule:
    \[
        \nabla(fs) = df \otimes s + f\nabla s,
    \]
    where $f \in C^\infty(M), s \in C^\infty(V)$.

    A connection tells us how to differentiate sections of the bundle $V$. Given a section $s \in C^\infty(V)$ and a 
    vector field $X \in C^\infty(TM)$, the section $\nabla_X s \in C^\infty(V)$ acts as the ``derivative''
    of $s$ in the direction of $X$. 
\end{definition}

\begin{example}
    An example of a connection on $\RR^n$ with $V = TM$ is the Lie derivative,
    \[
        \nabla_X: s \mapsto \mathcal{L}_X s
    \]
    where $X, s \in C^\infty(TM)$.
\end{example}

A related notion is \textbf{Parallel Transport}, which is a way to lift paths in $M$ to paths in 
our bundle $V$. 

\begin{definition}
    We define the \textbf{Curvature Operator} $K$ of $V$ to be the map
    \[
        K(X, Y)Z = \nabla_X \nabla_Y Z - \nabla_Y \nabla_X Z - \nabla_{[X,Y]}Z.
    \]
    In other words, $K(X,Y)$ is the difference $[\nabla_X, \nabla_Y] - \nabla_{[X, Y]}$.

    We can also express this as a tensor called the \textbf{Curvature 2-form} 
    \[
        K = \sum_{i < j} K(\partial_i, \partial_j)dx^i \wedge dx^j
    \]
    which takes values in $\text{End}(V)$. In other words, each $K(\partial_i, \partial_j)$ is an element
    of $\text{End}(V)$. 
\end{definition}

Now let $M$ be equipped with a Riemannian Metric $(\cdot, \cdot)$. 

\begin{definition}
    A \textbf{Symmetric} connection on $TM$ is one such that
    \[
        \nabla_X Y - \nabla_Y X = [X, Y].
    \]
    We say a connection is \textbf{Compatible with the metric} if for any three vector 
    fields $X, Y_1,$ and $Y_2$, 
    \[
        (\nabla_X Y_1, Y_2) + (Y_1, \nabla_X Y_2) = X \cdot (Y_1, Y_2).
    \]
\end{definition}

\begin{theorem}[Levi-Civita]
    A Riemannian manifold possesses a unique connection called the \textbf{Levi-Civita Connection} 
    that is symmetric and compatible with the metric.
\end{theorem}


\begin{definition}
    The curvature of the Levi-Civita connection on a riemannian manifold is called the 
    \textbf{Riemann Curvature Tensor} (or operator) and is denoted $R$. We write 
    \[
        R(e_j, e_k)e_l = \sum_{i}R^i_{ljk}e_i
    \]
    relative to some local frame $\{e_i\}$ for $TM$.
\end{definition}

\begin{proposition}
    The Riemann curvature operator has the following properties:
    \begin{enumerate}
        \item[(i)] R(X,Y)Z + R(Y,X)Z = 0
        \item[(ii)] R(X,Y)Z + R(Y,Z)X + R(Z,X)Y = 0
        \item[(iii)] (R(X,Y)Z, W) + (R(X,Y)W, Z) = 0
        \item[(iv)] (R(X,Y)Z, W) = (R(Z,W)X,Y).
    \end{enumerate}
\end{proposition}

These justify some calculations later.

\begin{definition}
    The \textbf{Ricci curvature tensor} is the bilinear form on $TM$ defined by 
    \[
        \text{Ric}(Y,Z) = \text{tr}(X \mapsto R(X,Y)Z),
    \]
    or $\text{Ric}_{ab} = \sum_i R_{aib}^i$.
\end{definition}

\begin{definition}
    We will define positive scalar curvature as the trace of the Ricci curvature operator, expressed as
    \[
        \kappa = \sum_{i,j}R_{ijij}
    \]
    in an orthonormal framing
\end{definition}

\section{Clifford algebras}

\begin{definition}
    Let $V$ be a vector space equipped with a symmetric bilinear form $(\cdot, \cdot)$. A 
    \textbf{Clifford algebra} for $V$, denoted $\text{Cl}(V)$ is a unital algebra $A$ which is equipped with a map 
    $\phi:V \to A$ such that $\phi(v)^2 = -(v,v)1$, which is universal among all such algebras. 

    It can also be viewed as the quotient of the tensor algebra $\mathcal{T}(V)$ of $V$ by the ideal 
    generated by elements of the form $v \otimes v +q(V)1$ for $v \in V$. 
\end{definition}

Here we will mostly take $V = \RR^k$, and write $\text{Cl}(k) = \text{Cl}(V)$. With a positive definite inner
product, there is only one group $\text{Cl}(k)$.

\begin{example}
    If we take the bilinear form on $V$ to be uniquely zero, we get $\text{Cl}(V) = \bigwedge^*(V)$ as 
    the clifford algebra. For $V = \RR^k$, we write $\text{Cl}(k)$.
\end{example}

\begin{definition}
    A \textbf{Clifford module} for the real inner product space $V = \RR^k$ is a left module over the complex algebra 
    $\mathbb{C}l(k) := \text{Cl}(k)\otimes_{\RR}\mathbb{C}$.

    This is the complexification of the algebra $Cl(V)$. It follows from the definition of the tensor product
    that for any $v \in \mathbb{C}l(V)$, we can write
    \[
        v = v_1 \otimes 1 + v_2 \otimes i,
    \]
    and thus there is a canonical direct sum decomposition
    \[
        \mathbb{C}l(V) = Cl(V) \otimes iCl(V).
    \]

    A \textbf{Clifford bundle} is a bundle of Clifford modules over a Riemannian manifold $M$ with 
    a compatible connection and metric, such that 
    \begin{enumerate}
        \item[(i)] The Clifford action of each $v \in T_m M$ on $S_m$ is skew-adjoint.
        \item[(ii)] The connection on $S$ is compatible with the Levi-Civita connection on $M$,
        in that it satisfies a product rule. 
    \end{enumerate}
\end{definition}

\begin{definition}
    There is a compact Lie group called $\text{Spin}(k)$ which is a double cover over $\text{SO}(k)$. For 
    $k \geq 3$, it is simply connected, and the universal covering space of $\text{SO}(k)$. It is a certain
    subgroup of $\text{Cl}(k)$ generated by some invertible elements. 

    Let $M$ be a manifold and $E$ a vector bundle on $M$. Let $P_{\text{SO}}(E)$ be the principal $\text{SO}(k)$
    bundle of frames for $E$. Then we define a \textbf{spin structure} on $E$ to be a principal $\text{Spin}(k)$-bundle
    $P_{\text{Spin}}(E)$ on $M$, which is a double cover of $P_{\text{SO}}(E)$, such that the global double
    cover is compatible with the fiberwise double cover of $\text{Spin}(k) \to \text{SO}(k)$. 

    A \textbf{spin manifold} is a manifold with a spin structure on its tangent bundle $TM$. 

    To every even dimensional spin manifold $M$, there is an associated vector bundle $S$ on $M$ called the 
    \textbf{spinor bundle} of dimension $2^n$, which is in fact a clifford bundle. There is a 
    canonical connection on $S$ called the spin connection, which comes from the Riemannian structure on $M$. 
\end{definition}

\section{Dirac operators}

In this section, let $M$ be a spin manifold of even dimension $n$, $S$ the spinor bundle, and $\nabla$ the 
spin connection on $S$. We will construct what is called the Atiyah-Singer operator $D$, which is the
dirac operator in the case of a spin manifold on its spinor bundle. 

\begin{definition}
    We define $D$ by the following sequence of maps 
    \[
        C^\infty(S) \to C^\infty(T^*M \otimes S) \to C^\infty(TM\otimes S) \to C^infty(S),
    \]
    where the first arrow is just the connection on $S$, the second arrow is the identification of $T^*M$
    with $TM$ using the metric on $M$, and the last arrow comes from the clifford action of sections of $TM$
    on $S$ fiberwise. 

    In a local orthonormal basis, we can write
    \[
        Ds = \sum_i e_i \nabla_i s
    \]
    where $\nabla_i$ is the connection applied to $\partial_i$. 

    The Dirac operator is an unbounded operator, and is (formally) self adjoint in the correct 
    function space, but we will not get into these details. 

    An important fact is that the spinor bundle is graded, and the dirac operator is chirality-reversing.
    Then we can write 
    \[
        D = \begin{bmatrix}
            0 & D^-\\
            D^+ & 0
        \end{bmatrix}.
    \]
    Moreover, these operators are adjoints of one another, and we have the important formula
    \[
        \text{Index}(D^+) = \dim\ker D^+ - \dim \ker D^-
    \]
    for the index of $D$.
\end{definition}

\begin{proposition}
    Let $e_i$ be an orthonormal frame at $m \in M$. Then we have
    \begin{align}
        D^2s & = \sum_{i,j}e_j \nabla_j(e_i \nabla_i s)\\
        & = \sum_{i,j}e_j e_i \nabla_j \nabla_i s\\
        & = -\sum_i \nabla_i^2 s + \sum_{j < i} e_j e_i (\nabla_j \nabla_i - \nabla_i \nabla_j)s\\
        & = \nabla^*\nabla s + Ks.
    \end{align}
    Here, the symbol $\nabla^* \nabla$ is used to indicate that $\nabla$ is a sort of differential
    operator from $S$ to $T^*M \otimes S$, and $\nabla^*$ is its adjoint. The negative sign is from
    the adjoint operator. The operator $K$ can be expanded
    \[
        K = F + \frac{1}{4}\kappa,
    \]
    where $\kappa$ is the scalar curvature of $M$ and $F$ is called the twisting curvature.
    The crucial fact about spin manifolds is that here, $F$ vanishes, leaving
    \[
        D^2 = \nabla^*\nabla + \frac{1}{4}\kappa,
    \]
    which is known as the \textbf{Lichnerowicz formula}.
\end{proposition}

An important construction is as follows. Let $M$ be a spin manifold and $S$ its spinor bundle,
and $E$ be any vector bundle on $M$. Then $S \otimes E$ is also a clifford bundle on $M$,
with a tensor product connection. We can define the dirac operator on this clifford bundle as well, and 
we will denote this operator by $D_E$. We have the analagous formula
\[
    D_E^2 = \nabla^*\nabla + \frac{1}{4}\kappa + R^E,
\]
where $R^E$ only depends on the curvature tensor of $R^E$ and the dimension, and $\kappa$ is the scalar curvature
of $M$. 

The term $R^E$ corresponds to how ``flat'' the bundle $E$ is. This is a ``Bochner formula''.

Now we state (without proof) a certain corollary of the Atiyah-Singer index theorem on spin manifolds:
\begin{theorem}
    The index of the positive part of the dirac operator is given by
    \[
        \text{Index}(D^+) = \hat{A}(X),
    \]
    where $\hat{A}(X)$ is a certain topological invariant of $X$. Moreover, if $E$ is a bundle and 
    $D_E$ is the twisted operator from before, then we have
    \[
        \text{Index}(D_E^+) = \{\text{ch}E \cdot \hat{\textbf{A}}(X)\}[X],
    \]
    which is another topological invariant of the bundle $E$.
\end{theorem}

\section{Final theorem}

\begin{definition}
    A map $f: X \to S^n$ is \textbf{$\epsilon$-contracting} if we have
    \[
        \|f^*v\| \leq \epsilon \|v\|
    \]
    for all tangent vectors $v$. 

    A compact riemannian manifold of dimension $n$ is \textbf{compactly enlargeable} if for every $\epsilon > 0$, there exists an orientable
    Riemann covering space which admits an $\epsilon$-contracting map to $S^n$ of non-zero degree, such that
    the covering map finite.
\end{definition}

An example of an enlargeable manifold is the $n$-torus $\mathbb{T}^n$. I can draw the map.

Now here is the final theorem:

\begin{theorem}
    An enlargeable spin manifold may not carry a metric of positive scalar curvature. 
\end{theorem}

\textbf{Proof:} Let $X$ be an even dimensional compact spin manifold of even dimension. The 
case for odd dimension will be mentioned at the end. Suppose that $X$ does carry a metric 
of positive scalar curvature, with $\kappa \geq \kappa_0 + c > 0$ for some $\kappa_0$ and $c > 0$. 

Choose a complex vector bundle $E_0$ on the sphere $S^n$ such that the top Chern class $c_n(E_0) \neq 0$. 
This is always possible. Let $\nabla^{E_0}$ be a unitary connection.

Let $\epsilon > 0$ and choose a finite orientable covering $\tilde{X} \to X$ which admits 
an $\epsilon$-contracting map to the sphere of non-zero degree. Use $f$ to pull back $E_0$
to a bundle $E$ on $\tilde{X}$. Consider the Atiyah-Singer operator $D_E$ on $S\otimes E$ 
as mentioned before. We apply the Bochner formula
\[
    D_E^2 = \nabla^*\nabla + \frac{1}{4}\kappa + R^E,
\]
where $R^E$ again depends only on the curvature of the bundle $E$. This operator is a bundle homomorphism.

We will now use (without proof) that the term $R^E$ satisfies the estimate
\[
    \|R^E\| < \frac{1}{4}k_0
\]
on each fiber of $M$. Let $s$ be a section of $S\otimes E$. Then pointwise, we have
\begin{align}
    \langle D_E^2 s, s\rangle & = |\nabla s|^2 + \frac{1}{4}\kappa |s|^2 + \langle R^E s, s\rangle.
\end{align}
For the last term, we have
\[
    \langle R^E s, s \rangle \leq \|R^E\||s|^2 < \frac{1}{4}k_0 |s|^2,
\]
so that
\[
    \langle D_E^2 s, s\rangle > |\nabla s|^2 + \frac{1}{4}\kappa |s|^2 - \frac{1}{4}|s|^2 > |\nabla s|^2 + c|s|^2.
\]

This gives us the global estimate
\[
    \|D_E^2 s\| \geq c\|s\|.
\]
Now it is clear that if $D_E s = 0$, we must have $s = 0$. In fact this implies $D_E^+ s = 0$, so that 
$\ker D^+ = \ker D^- = \{0\}$. In particular, we have
\[
    \text{Index}(D_E^+) = \dim\ker D^+ - \dim\ker D^- = 0.
\]

However, from a calculation using the Index formula for twisted dirac operators, we can show the index of 
the dirac operator is non-zero. This gives a contradiction, and thus $X$ may not admit a metric 
of positive scalar curvature. 

\end{document}

